Dans la suite, nous considérons l'ensemble $\setproba{M}$ des matrices $A$ telles que $A^T Q A = Q$ et $\det A = 1$ afin de pouvoir les multiplier entre elles comme nous allons le vérifier.%
\footnote{
    Ce qui suit reprend les excellentes indications données par Jérôme Germoni dans une discussion au bas de cette page:
\url{http://images.math.cnrs.fr/+Nombres-puissants-au-bac-S+}
\textit{(chercher les messages de l'utilisateur projetmbc)}.
}


\medskip

Nous savons que $A \in \setproba{M}$ est du type
$A
=
\begin{pmatrix}
  a & 8c \\
  c & a
\end{pmatrix}$
et que la contrainte $a^2 - 8c^2 = 1$ est cachée dans la relation $A^T Q A = Q$.%
\footnote{
    Autrement dit, la première colonne de $A$ est une solution de \ED.
}


\medskip

Nous savons aussi que si
$\begin{pmatrix}
  x \\
  y
\end{pmatrix}$
verifie $x^2 - 8 y^2 = 1$, il en sera de même pour
$\begin{pmatrix}
  x' \\
  y'
\end{pmatrix}
=
A
\begin{pmatrix}
  x \\
  y
\end{pmatrix}$,
c'est-à-dire
$\begin{pmatrix}
  x' \\
  y'
\end{pmatrix}
=
\begin{pmatrix}
  a x + 8c y \\
  c x + a y
\end{pmatrix}$.


\bigskip

Notant $\setproba{S}$ l'ensemble des solutions entières
\footnote{
	Ce qui suit se généralise aux solutions rationnelles, à celles réelles, ou encore à celles complexes, mais ne sortons pas du cadre de ce modeste document.
}
de \ED, ce qui précède motive la définition de la loi $\star$ sur $\setproba{S}$ via
$\begin{pmatrix}
  a \\
  c
\end{pmatrix}
\star
\begin{pmatrix}
  x \\
  y
\end{pmatrix}
=
\begin{pmatrix}
  a x + 8c y \\
  c x + a y
\end{pmatrix}$.
Nous allons démontrer que cette loi $\star$ est bien définie et fait de $\setproba{S}$ un groupe. Très joli ! Non?


\medskip

Les justifications suivantes sont faciles à obtenir en notant que
$\begin{pmatrix}
  a \\
  c
\end{pmatrix}
\star
\begin{pmatrix}
  x \\
  y
\end{pmatrix}$
est la première colonne de la matrice produit
$\begin{pmatrix}
  a & 8c \\
  c & a
\end{pmatrix}
\begin{pmatrix}
  x & 8y \\
  y & x
\end{pmatrix}$.
\footnote{
	Les résultats présentés ici se vérifient et se trouvent directement à la main sans passer par les matrices, mais ceci est bien moins élégant.
}


\begin{itemize}
	\item La loi $\star$ est bien interne à $\setproba{S}$, car si $A^T Q A = Q$ avec $\det A = 1$ et $B^T Q B = Q$ avec $\det B = 1$, nous avons:

	\begin{itemize}
		\item $(AB)^T Q (AB) = B^T A^T Q A B = B^T Q B = Q$

		\medskip
		\item $\det(AB) = \det A \cdot \det B = 1$
	\end{itemize}


	\medskip
	\item
	$\begin{pmatrix}
	  1 \\
	  0
	\end{pmatrix}$
	est l'élément neutre de $\star$, car sa matrice associée est la matrice identité.


	\medskip
	\item
	$\begin{pmatrix}
	  a  \\
	  -c
	\end{pmatrix}$
	est l'inverse de
	$\begin{pmatrix}
	  a \\
	  c
	\end{pmatrix}$,
	car
	$\begin{pmatrix}
	  a  & -8c \\
	  -c & a
	\end{pmatrix}$
	est l'inverse matriciel de
	$\begin{pmatrix}
	  a & 8c \\
	  c & a
	\end{pmatrix}$
    (voir aussi le point précédent).


	\medskip
	\item Pour finir, l'associativité de $\star$ découle de celle de la multiplication des matrices.
\end{itemize}
